\section{A source-restricted odd deletion with exceptional minima}
\label{sec:restricted-odd-capacity}

Let $D$ be strongly connected, all-deficient and minimal under deletion
of arbitrary nonempty sets of arcs, with $n=2\delta+3$.
Use $e(v)=d_D^+(v)-\delta$, $b(v)=g(v)-1\ge0$, and the heavy-target
notation established earlier.

\begin{lemma}[A three-vertex deletion with $d$ exceptional sources]
\label{lem:restricted-odd-degree}
Let $S$ consist of three heavy targets.  Suppose every surviving
minimum-degree vertex sends at most two arcs into $S$, and every
surviving nonminimum-degree vertex has outdegree at least $\delta-1$
after deletion of $S$.  Let
\[
 Z=\{v\notin S:e(v)=0,\ |N_D^+(v)\cap S|=2\},\qquad d=|Z|.
\]
Then
\begin{equation}\label{eq:restricted-odd-degree-bound}
 {\delta\le20d+12.}
\end{equation}
If $\delta\ge10$, then $d\ge1$.
\end{lemma}

\begin{proof}
Put $D'=D-S$, $n'=2\delta$, and fix the counting baseline
$\alpha=\delta-1$.  Set $W=V(D')\setminus Z$.
The vertices of $Z$ have outdegree $\alpha-1$, while every vertex
of $W$ has outdegree at least $\alpha$.
When $d>0$, the actual minimum outdegree of $D'$ is $\delta-2$;
we do not apply a full-graph minimum-degree theorem at baseline
$\alpha$.  All capacity arguments below use only the source set $W$.

Write $a'(v)=d_{D'}^+(v)-\alpha$, let $\mu'(v)$ be the missing
degree in $D'$, and define
\[
 T'(x)=V(D')\setminus N^-_{D'}[x],\qquad
 q'(x)=|T'(x)|-\alpha=a'(x)+\mu'(x),
\]
\[
 P'_W(x)=\{v\in W:x\in U_{D'}(v)\},\quad
 \mathcal X=\{x:P'_W(x)\ne\varnothing\},\quad N=|\mathcal X|.
\]
For $x\in\mathcal X$ put
$p_x=|P'_W(x)|$ and $\Delta_x=2q'(x)-1-p_x$.
The source-restricted capacity count gives $\Delta_x\ge0$.
Let $\mathcal D=\sum_{x\in\mathcal X}\Delta_x$.
The only negative $q'$ values are $q'=-1$ at
\[
 Z_0=\{z\in Z:\mu'(z)=0\},\qquad z_0=|Z_0|,
\]
and these vertices are never targets.
Define the nonnegative quantities
\[
 t_2=2\!\sum_{x\notin\mathcal X,\ x\notin Z_0}\!q'(x),
 \qquad t_3=\sum_{v\in W}b(v),
\]
\[
 E_S(v)=(N_D^{++}(v)\setminus S)\setminus N_{D'}^{++}(v),
 \qquad t_4=\sum_{v\in W}|E_S(v)|.
\]
Let $\Phi(S)$ be the exact deletion imbalance of
Lemma~\ref{lem:odd-root-transfer}, and put
$Q_Z=\sum_{v\in Z}|S\cap N_D^{++}(v)|$.
Each vertex of $Z$ directly dominates two of the three deleted
vertices, so $0\le Q_Z\le d$.

The exact restricted budget is
\begin{equation}\label{eq:restricted-odd-budget}
 {N+\mathcal D+t_2+t_3+t_4
                =\beta:=2d-\Phi(S)+2z_0+Q_Z.}
\end{equation}
To derive it, put $A'=\sum_{v\in V(D')}a'(v)$ and
$B'_W=\sum_{v\in W}b_{D'}(v)$.
The even-order arc count gives $\sum_xq'(x)=n'-A'$.
Since $a'=-1$ on $Z$, double counting the far source--target pairs
with sources in $W$ gives
\[
 N+\mathcal D+2\sum_{x\notin\mathcal X}q'(x)+B'_W=4d.
\]
The exact transfer on $W$ is
\[
 B'_W=t_3+\Phi(S)-Q_Z+2d+t_4.
\]
Move the $-2z_0$ nontarget contribution to the other side to obtain
\eqref{eq:restricted-odd-budget}.
For the heavy triple, $\Phi(S)\ge-3$ by its exact internal-pair
formula.  Hence
\begin{equation}\label{eq:restricted-odd-beta-coarse}
 \beta\le5d+3.
\end{equation}

Partition the nonexceptional sources into
\[
 \Pi=\{v\in W:q'(v)=0\},\qquad Y=W\setminus\Pi,
 \qquad y=|Y|.
\]
Vertices of $\Pi$ are whole of outdegree $\alpha$ in $D'$;
all missing pairs of $D'$ lie inside $Y\cup Z$.
Moreover $y\le N+t_2/2$.
The integer light-target charging argument of
Section~\ref{sec:predecessor-bound}, applied with sources $W$ and
baseline $\alpha$, gives
\begin{equation}\label{eq:restricted-odd-charge}
 M(D')\le\mathcal D+\frac{t_2}{2}
                   +\sum_{v\in Y}|U_{D'}(v)|+d.
\end{equation}
The additional $d$ accounts for the term $[x\in Z]$ in
$\mu'(x)\le q'(x)+[x\in Z]$.
The vertices of $Z_0$ have missing degree zero and incur no
missing-pair charge.
No full-graph capacity assertion is used in this charging step.

Write
\[
 A_Y=\sum_{v\in Y}e(v),\quad
 J_Y=\sum_{v\in Y}|N_D^+(v)\cap S|,\quad
 A'_Y=\sum_{v\in Y}a'(v).
\]
Then $A'_Y=A_Y+y-J_Y$.
The exact degree count in $D'$ gives
\[
 M(D')=\delta+d-A'_Y.
\]
The far-set transfer, retaining the original excess, gives
\[
 \sum_{v\in Y}|U_{D'}(v)|
 =3y-2A_Y+\sum_Yb(v)-\sum_Y|S\cap U_D(v)|
                       +\sum_Y|E_S(v)|
 \le3y-2A_Y+t_3+t_4.
\]
Substitution into \eqref{eq:restricted-odd-charge} yields
\[
 \delta+A_Y+J_Y
 \le\mathcal D+\frac{t_2}{2}+4y+t_3+t_4
 \le4N+\mathcal D+\frac{5t_2}{2}+t_3+t_4
 \le4\beta.
\]
In fact the retained version is
\[
 \delta+A_Y+J_Y
 \le4\beta-3\mathcal D-\frac{3t_2}{2}-3t_3-3t_4.
\]
Dropping the nonnegative terms and using
\eqref{eq:restricted-odd-beta-coarse} proves
\eqref{eq:restricted-odd-degree-bound}.
If $d=0$, the deletion is admissible for the heavy-triple
obstruction, which rules out $\delta\ge10$.
\end{proof}

\begin{theorem}[A residual-dependent restriction on four heavy roots]
\label{thm:four-heavy-residual}
Assume $\delta\ge10$ and $H=4$.
Let $m$ count minimum-degree vertices in $K$, let $L$ count
nonwhole vertices outside $K$, and put $E=m+2L$.
Then $E\ge1$ and
\begin{equation}\label{eq:four-heavy-degree-residual}
 {\delta\le20(m+2L)+12\le40R+12.}
\end{equation}
Consequently
\[
 M\ge\left\lceil\frac{81\delta+108}{80}\right\rceil
 \qquad\text{when }H=4.
\]
\end{theorem}

\begin{proof}
For a heavy root, the general predecessor bound gives
$15c\ge\delta+4$.
If some heavy root has $c\le E$, then
$\delta\le15E-4\le20E+12$.

Otherwise every heavy root has more than $E$ minimum-degree
predecessors.  Excluding the at most $E$ bad witnesses described
in Theorem~\ref{thm:general-good-witnesses} leaves good whole
minimum-degree predecessors at every root.
Use the \emph{full} oriented relation witnessed by all such good
sources: $i\to j$ if a good source points to $r_i$ and is far
from $r_j$.  It has positive weighted outdegree at every root
and no digon, since both witnesses of a putative digon would
be whole.

By Lemma~\ref{lem:odd-root-three-disjoint}, choose three roots $S$
whose outneighborhoods in this relation are pairwise disjoint.
A good minimum-degree source cannot point to two chosen roots:
its nonempty far certified core set would place some far root
in both outneighborhoods.  Thus only bad minimum sources can
lose two arcs into $S$, and their number $d$ is at most $E$.
The universal heavy-capacity bound at $H=4$ limits every vertex
to two outgoing arcs into the full heavy set.  All other
minimum sources lose at most one arc, and nonminimum sources
retain outdegree at least $\delta-1$.
Lemma~\ref{lem:restricted-odd-degree} therefore gives
$\delta\le20d+12\le20E+12$.
For $\delta\ge10$ that lemma also gives $d\ge1$.
If $E=0$, the first case is impossible and the second would have
$d=0$, so $E\ge1$.

Finally $2m+L\le R$ implies $m+2L\le2R$.
This proves \eqref{eq:four-heavy-degree-residual}.
Substitute $R=2M-(2\delta+3)$ and rearrange to obtain
$80M\ge81\delta+108$, giving the stated integer bound.
\end{proof}
